% ID: SEL-D0034-OA-E02
% ESTADO: validada. S035, cotejada con examen y criterios.
\textbf{a)} La parábola de $f'$ abre hacia arriba, porque
su vértice está por debajo de las dos raíces. Su signo es
positivo en $(-\infty,-1)$ y $(3,\infty)$, y negativo en $(-1,3)$.
Por tanto, $f$ crece hasta $-1$, decrece entre $-1$ y $3$
y crece a partir de $3$. Hay un máximo relativo en $x=-1$
y un mínimo relativo en $x=3$, por los cambios de signo.

\textbf{b)} Del vértice de la gráfica de $f'$ se lee $f'(1)=-2$.
Como $f(1)=2$, la tangente es $y-2=-2(x-1)$,
es decir, $\boxed{y=4-2x}$.
